\section{Lineární regrese}\label{sec:ai-linearni-regrese}

Lineární regrese hledá funkci
\[
    f(x)=ax+b,
\]
která co nejlépe popisuje vztah mezi statistickými soubory
\[
    \mathbf{X}=(x_1,\ldots,x_n),
    \qquad
    \mathbf{Y}=(y_1,\ldots,y_n).
\]
Hodnota
\[
    \hat y_i=f(x_i)
\]
je předpověď modelu pro vstup \(x_i\), zatímco
\[
    e_i=y_i-\hat y_i
\]
nazýváme \emph{reziduum}. Kladné reziduum znamená, že bod leží nad regresní přímkou, záporné reziduum že leží pod ní.

\begin{figure}[ht]
    \centering
    \begin{tikzpicture}
    \begin{axis}[
        width=12cm,height=7cm,
        xmin=0,xmax=7,ymin=0,ymax=6,
        axis lines=left,xlabel={$x$},ylabel={$y$},
        xtick=\empty,ytick=\empty,
        label style={font=\normalsize},
        every node/.style={font=\normalsize},
        clip=false
    ]
        \addplot[red!75!black,very thick,domain=0.5:6.5]{0.7*x+0.5};
        \addplot[only marks,mark=*,mark size=2.5pt,blue!75!black]
        coordinates{(1,1.8)(2,1.5)(3,3.2)(4,3.0)(5,4.7)(6,4.5)};
        \draw[dashed,gray] (axis cs:1,1.8) -- (axis cs:1,1.2);
        \draw[dashed,gray] (axis cs:2,1.5) -- (axis cs:2,1.9);
        \draw[dashed,gray] (axis cs:3,3.2) -- (axis cs:3,2.6);
        \draw[dashed,gray] (axis cs:4,3.0) -- (axis cs:4,3.3);
        \draw[dashed,gray] (axis cs:5,4.7) -- (axis cs:5,4.0);
        \draw[dashed,gray] (axis cs:6,4.5) -- (axis cs:6,4.7);
        \node[red!75!black] at (axis cs:5.4,5.1){\(f(x)=ax+b\)};
        \node[gray] at (axis cs:3.35,2.7){\(e_i\)};
    \end{axis}
\end{tikzpicture}

    \caption{Předpovědi regresní přímky a~svislá rezidua.}
    \label{fig:ai-regrese-rezidua}
\end{figure}

\begin{definition}\label{def:ai-sse-mse}
    Nechť \(\mapping{f}{\R}{\R}\). \emph{Součet čtvercových chyb} a~\emph{střední kvadratickou chybu} funkce \(f\) definujeme jako
    \[
        \SSE_{\mathbf{X},\mathbf{Y}}(f)
        =
        \sum_{i=1}^{n}
        \bigl(y_i-f(x_i)\bigr)^2
    \]
    a
    \[
        \MSE_{\mathbf{X},\mathbf{Y}}(f)
        =
        \frac1n
        \SSE_{\mathbf{X},\mathbf{Y}}(f).
    \]
\end{definition}

Pouhý součet reziduí není vhodný, protože kladné a~záporné chyby se mohou vyrušit. Čtverce jsou vždy nezáporné a~větší chyby navíc zvýrazní. Minimalizace \(\SSE\) a~\(\MSE\) vede ke stejné přímce, protože se obě veličiny liší pouze kladným násobkem \(1/n\).

\subsection{Metoda nejmenších čtverců}

Pro model \(f_{a,b}(x)=ax+b\) hledáme
\[
    (a,b)
    =
    \argmin_{(a,b)\in\R^2}
    \sum_{i=1}^{n}(y_i-ax_i-b)^2.
\]
Nejprve zvolme pevné \(a\). Výraz lze chápat jako součet čtvercových odchylek hodnot statistického souboru
\[
    \mathbf{Y}-a\mathbf{X}
    =
    (y_1-ax_1,\ldots,y_n-ax_n)
\]
od konstanty \(b\). Podle
\smartref[tvrzení~][]{\showrefs}{prop:ai-prumer-vlastnosti}{minimalizační vlastnosti průměru} je nejlepší volbou
\[
    b
    =
    \average{\mathbf{Y}-a\mathbf{X}}
    =
    \average{\mathbf{Y}}
    -
    a\average{\mathbf{X}}.
\]
Každý zbývající kandidát proto prochází bodem
\[
    \bigl(\average{\mathbf{X}},\average{\mathbf{Y}}\bigr)
\]
a~má tvar
\[
    f(x)
    =
    \average{\mathbf{Y}}
    +
    a(x-\average{\mathbf{X}}).
\]

\begin{proposition}\label{prop:ai-regresni-primka}
    Nechť \(\Var(\mathbf{X})>0\). Potom existuje právě jedna lineární funkce
    \(f(x)=ax+b\), která minimalizuje
    \(\SSE_{\mathbf{X},\mathbf{Y}}\), a~její koeficienty jsou
    \[
        a
        =
        \frac{\Cov(\mathbf{X},\mathbf{Y})}
        {\Var(\mathbf{X})},
        \qquad
        b
        =
        \average{\mathbf{Y}}
        -
        a\average{\mathbf{X}}.
    \]
\end{proposition}

\begin{proof}
    Z~předchozí úvahy stačí minimalizovat chybu přímek procházejících bodem průměrů. Položme
    \[
        u_i=x_i-\average{\mathbf{X}},
        \qquad
        v_i=y_i-\average{\mathbf{Y}}.
    \]
    Po dosazení \(b=\average{\mathbf{Y}}-a\average{\mathbf{X}}\) dostaneme
    \begin{align*}
        \SSE_{\mathbf{X},\mathbf{Y}}(f)
        &=
        \sum_{i=1}^{n}(v_i-au_i)^2\\
        &=
        \sum_{i=1}^{n}v_i^2
        -2a\sum_{i=1}^{n}u_iv_i
        +a^2\sum_{i=1}^{n}u_i^2.
    \end{align*}
    Označme
    \[
        a^\ast
        =
        \frac{\sum_{i=1}^{n}u_iv_i}
        {\sum_{i=1}^{n}u_i^2}.
    \]
    Jmenovatel je kladný, protože \(\Var(\mathbf{X})>0\). Doplněním na čtverec získáme
    \begin{align*}
        \SSE_{\mathbf{X},\mathbf{Y}}(f)
        &=
        \sum_{i=1}^{n}v_i^2
        -
        \frac{\left(\sum_{i=1}^{n}u_iv_i\right)^2}
        {\sum_{i=1}^{n}u_i^2}\\
        &\quad+
        \left(\sum_{i=1}^{n}u_i^2\right)
        (a-a^\ast)^2.
    \end{align*}
    Poslední člen je nezáporný a~je nulový právě pro \(a=a^\ast\). Minimum je proto jediné a
    \[
        a^\ast
        =
        \frac{
            \frac1n\sum_{i=1}^{n}
            (x_i-\average{\mathbf{X}})
            (y_i-\average{\mathbf{Y}})
        }{
            \frac1n\sum_{i=1}^{n}
            (x_i-\average{\mathbf{X}})^2
        }
        =
        \frac{\Cov(\mathbf{X},\mathbf{Y})}
        {\Var(\mathbf{X})}.
    \]
    Jednoznačnost \(b\) pak plyne ze vztahu
    \(b=\average{\mathbf{Y}}-a\average{\mathbf{X}}\).
\end{proof}

Pomocí korelace lze směrnici zapsat také
\[
    a
    =
    \rho_{\mathbf{X}\mathbf{Y}}
    \frac{\sigma_{\mathbf{Y}}}{\sigma_{\mathbf{X}}},
\]
a~regresní přímku ve tvaru
\[
    f(x)
    =
    \average{\mathbf{Y}}
    +
    \rho_{\mathbf{X}\mathbf{Y}}
    \frac{\sigma_{\mathbf{Y}}}{\sigma_{\mathbf{X}}}
    (x-\average{\mathbf{X}}).
\]
Znaménko směrnice je tedy stejné jako znaménko korelace.

Uvedené vztahy poskytují přímo algoritmus pro trénování jednoduché lineární regrese. Není třeba zkoušet různé přímky ani opakovaně upravovat parametry: stačí jedním průchodem dat spočítat potřebné součty, z~nich průměry, rozptyl a~kovarianci a~nakonec koeficienty \(a,b\).

\begin{algorithm}
    \caption{\textsc{linear-regression}}
    \label{alg:ai-linear-regression}
    \KwIn{Statistické soubory \(\mathbf{X}=(x_1,\ldots,x_n)\) a~\(\mathbf{Y}=(y_1,\ldots,y_n)\) s~\(\Var(\mathbf{X})>0\)}
    \(\displaystyle\average{\mathbf{X}}\gets\frac1n\sum_{i=1}^{n}x_i\), \(\displaystyle\average{\mathbf{Y}}\gets\frac1n\sum_{i=1}^{n}y_i\)\\
    \(\displaystyle\Var(\mathbf{X})\gets\frac1n\sum_{i=1}^{n}(x_i-\average{\mathbf{X}})^2\)\\
    \(\displaystyle\Cov(\mathbf{X},\mathbf{Y})\gets\frac1n\sum_{i=1}^{n}(x_i-\average{\mathbf{X}})(y_i-\average{\mathbf{Y}})\)\\
    \(\displaystyle a\gets\frac{\Cov(\mathbf{X},\mathbf{Y})}{\Var(\mathbf{X})}\)\\
    \(b\gets\average{\mathbf{Y}}-a\average{\mathbf{X}}\)\\
    \KwOut{Regresní přímka \(f(x)=ax+b\)}
\end{algorithm}

\subsection{Minimální MSE regresní přímky}

\begin{proposition}\label{prop:ai-regrese-min-mse}
    Nechť \(\sigma_{\mathbf{X}}>0\), \(\sigma_{\mathbf{Y}}>0\) a~\(f\) je regresní přímka souborů \(\mathbf{X},\mathbf{Y}\). Potom
    \[
        \MSE_{\mathbf{X},\mathbf{Y}}^{\min}
        =
        \min_{\ell\ \text{lineární}}
        \MSE_{\mathbf{X},\mathbf{Y}}(\ell)
        =
        \MSE_{\mathbf{X},\mathbf{Y}}(f)
        =
        \sigma_{\mathbf{Y}}^2
        \left(1-\rho_{\mathbf{X}\mathbf{Y}}^2\right).
    \]
\end{proposition}

\begin{proof}
    Použijeme označení \(u_i,v_i\) z~předchozího důkazu. Pro optimální směrnici
    \[
        a=
        \frac{\sum_{i=1}^{n}u_iv_i}
        {\sum_{i=1}^{n}u_i^2}
    \]
    platí
    \begin{align*}
        \SSE_{\mathbf{X},\mathbf{Y}}(f)
        &=
        \sum_{i=1}^{n}(v_i-au_i)^2\\
        &=
        \sum_{i=1}^{n}v_i^2
        -2a\sum_{i=1}^{n}u_iv_i
        +a^2\sum_{i=1}^{n}u_i^2\\
        &=
        \sum_{i=1}^{n}v_i^2
        -
        \frac{
            \left(\sum_{i=1}^{n}u_iv_i\right)^2
        }{
            \sum_{i=1}^{n}u_i^2
        }.
    \end{align*}
    Z~definice korelace máme
    \[
        \rho_{\mathbf{X}\mathbf{Y}}^2
        =
        \frac{
            \left(\sum_{i=1}^{n}u_iv_i\right)^2
        }{
            \left(\sum_{i=1}^{n}u_i^2\right)
            \left(\sum_{i=1}^{n}v_i^2\right)
        }.
    \]
    Druhý člen předchozího výrazu je proto
    \[
        \frac{
            \left(\sum_{i=1}^{n}u_iv_i\right)^2
        }{
            \sum_{i=1}^{n}u_i^2
        }
        =
        \rho_{\mathbf{X}\mathbf{Y}}^2
        \sum_{i=1}^{n}v_i^2.
    \]
    Dostáváme
    \[
        \SSE_{\mathbf{X},\mathbf{Y}}(f)
        =
        \left(1-\rho_{\mathbf{X}\mathbf{Y}}^2\right)
        \sum_{i=1}^{n}v_i^2.
    \]
    Po vydělení \(n\) a~použití
    \[
        \frac1n\sum_{i=1}^{n}v_i^2
        =
        \Var(\mathbf{Y})
        =
        \sigma_{\mathbf{Y}}^2
    \]
    získáme požadovaný vztah.
\end{proof}

\subsection{Implementace v~Pythonu}

Program~\ref{lst:ai-linear-regression} nejprve spočítá průměry, rozptyl a~kovarianci a~potom přímo použije vzorce z~
\smartref[tvrzení~][]{\showrefs}{prop:ai-regresni-primka}{tvrzení o~regresní přímce}.

Metoda \texttt{fit} kontroluje také předpoklad \(\Var(\mathbf{X})>0\). Kdyby byly všechny vstupy stejné, data by ležela na jediné svislé přímce a~směrnici funkce \(y=ax+b\) by z~nich nebylo možné jednoznačně určit. Metoda proto místo dělení nulou oznámí chybu.

\lstinputlisting[
    caption={Jednoduchá lineární regrese metodou nejmenších čtverců.},
    label={lst:ai-linear-regression}
]{07-ai/code/linear_regression.py}

Trénování vyžaduje pouze konečný počet průchodů seznamy \texttt{x} a~\texttt{y}, takže má časovou složitost \(\bigO(n)\) a~kromě vstupních dat spotřebuje konstantní pomocnou paměť. Po natrénování je jedna předpověď jen výpočtem \(ax+b\), tedy operací v~konstantním čase.
